% cpu/hwfreelunch.tex
% mainfile: ../perfbook.tex
% SPDX-License-Identifier: CC-BY-SA-3.0

\section{硬件的免费午餐？}
\label{sec:cpu:Hardware Free Lunch?}
%
\epigraph{当今最大的问题是有太多人在寻找别人来为自己做些什么。
	  大多数问题的解决之道在于每个人为自己做些什么。}
	 {Henry Ford，改编}

近年来并发性受到如此多关注的主要原因是
\IXaltr{Moore's-Law}{Moore's Law}驱动的单线程性能提升
（或曰``免费午餐''~\cite{HerbSutter2008EffectiveConcurrency}）的终结，
如\cref{fig:intro:Clock-Frequency Trend for Intel CPUs}
（第~\cpageref{fig:intro:Clock-Frequency Trend for Intel CPUs}页）所示。
本节简要介绍硬件设计师可能采用的一些办法，
这些办法可以带来某种形式上的``免费午餐''。

然而，前一节展示了利用并发性所面临的一些重大硬件障碍。
其中对硬件设计师来说最严重的一个限制莫过于有限的光速。
如\cref{fig:cpu:System Hardware Architecture}
（第~\cpageref{fig:cpu:System Hardware Architecture}页）所述，
在1.8\,GHz的时钟周期内，光在真空中只能传播约8厘米远。
对于5\,GHz时钟，这个距离更是降到约3厘米。
这两个距离对于一个现代计算机系统的尺寸来说，还是太小了点。

更糟糕的是，电子在硅中的移动速度是真空中光速的1/30到1/3，
而常见的时钟逻辑结构运行得更慢。例如，
一次内存引用可能需要等待本地cache查找完成后才能将请求传递给系统的其余部分。
此外，将电信号从一个硅芯片传输到另一个（例如CPU与主存之间的通信）
需要相对低速且高功耗的驱动器。

\QuickQuiz{
	但单个电子的移动速度远没有那么快，即使在导体中也是如此！
	在半导体电压水平下，导体中电子的漂移速度仅为每秒
	约一\emph{毫米}。
	这是怎么回事？
}\QuickQuizAnswer{
	电子漂移速度追踪的是单个电子的长期运动。
	事实上，单个电子的运动相当随机，
	因此它们的瞬时速度很高，
	但从长期来看，它们并没有移动多远。
	在这一点上，电子类似于长途通勤者，
	他们可能大部分时间以全高速公路速度行驶，
	但从长期来看却哪里也没去。
	这些通勤者的速度可能是每小时70英里（每小时113公里），
	但他们相对于地球表面的长期漂移速度是零。

	因此，我们应该关注的不是电子的漂移速度，
	而是它们的瞬时速度。
	然而，即使是它们的瞬时速度也远不及光速的很大一部分。
	尽管如此，导体中电波的实测速度\emph{确实}是光速的相当大一部分，
	所以我们手头仍有一个谜。

	另一个诀窍是，电子之间在相当大的距离上
	（至少从原子的角度来看）相互作用，
	这要归功于它们的负电荷。
	这种相互作用由光子承载，而光子\emph{确实}以光速运动。
	所以即使在电学中的电子那里，
	大部分高速工作也是由光子完成的。

	推广通勤者的类比，一个司机可能使用智能手机
	通知其他司机有关事故或拥堵的情况，
	从而使交通流量的变化传播速度远快于
	单个汽车的瞬时速度。
	总结电学与交通流的类比：

	\begin{enumerate}
	\item	电子的（非常低的）漂移速度类似于通勤者的长期速度，
		两者都非常接近于零。
	\item	电子的（仍然相当低的）瞬时速度类似于
		交通中汽车的瞬时速度。
		两者都远高于漂移速度，
		但与变化传播速度相比仍然很小。
	\item	电波的（高得多的）传播速度主要是由于光子
		在电子之间传递电磁力。
		同样，由于司机之间的通信，
		交通模式可以相当快地改变。
		当然，这对已经堵在路上的司机不一定有多大帮助，
		就像对已经聚集在某个电容器中的电子一样。
	\end{enumerate}

	要完全理解这个主题，请学习电动力学。
}\QuickQuizEnd

事实上，据说Stephen Hawking曾声称半导体制造商面临两个基本问题：
\begin{enumerate*}[(1)]
\item 有限的光速，以及
\item 物质的原子性质~\cite{BryanGardiner2007}。
\end{enumerate*}
没错，光太慢而原子太大！

尽管如此，仍有一些技术（包括硬件和软件）可能有助于改善现状：

\begin{enumerate}
\item	新型材料和工艺，
\item	用光代替电，
\item	3D集成，
\item	专用加速器，以及
\item	现有的并行软件。
\end{enumerate}

以下各节分别介绍上述每一项。

\subsection{新型材料和工艺}
\label{sec:cpu:Novel Materials and Processes}

半导体制造商很有可能已经逼近了有限光速和
非零原子尺寸所隐含的极限。
然而，仍有一些研究和开发方向致力于绕过这些基本物理定律。

其中一个规避物质原子性质的办法是一种称为``高K介电''的材料，
这种材料允许较大的器件模拟不可行的微小器件的电学特性。
这些材料存在一些严峻的制造困难，
但仍然可能有助于将前沿推得更远一些。
另一种比较奇特的变通方法是在单个电子中存储多个比特，
利用给定电子可以存在于多个能量级别的事实。
这种特定方法能否在量产半导体器件中稳定工作还有待观察。

另一种提出的变通方法是``量子点''方法，
它允许更小的器件尺寸，但仍处于研究阶段。
还有一种提出的变通方法是用真空代替传统晶体管底部的原子，
从而得到真空间隙晶体管。

一个挑战是，许多最近的硬件器件级突破
需要非常精确地控制哪个原子放在
哪里~\cite{MichaelJKelly2017DeviceLevel}。
因此，找到在芯片上数十亿个器件中逐个放置原子的好方法的人，
即使别无所获，至少也能获得最出色的吹牛资本！

\subsection{光，而非电子}
\label{sec:cpu:Light; Not Electrons}

虽然光速是一个很难逾越的限制，但半导体设备更多是受限于电子移动的速度，
而非光速本身。半导体材料中电波的传播速度仅为
真空光速的3\pct 到30\pct。
在硅质器件上使用铜导线是一种提升电传播速度的方法，
而且很有可能进一步的进展将让电传播速度更加接近实际光速。
此外，已经有一些实验使用微型光纤作为芯片内部和芯片之间的传输通道，
因为在玻璃中光速能达到真空光速的60\pct 甚至更多。
这种方法存在的一个问题是光电/电光转换的效率，
会产生功耗和散热问题。

不仅如此，在物理学领域没有突破性进展的情况下，
数据流传输速度的任何指数增长都将受到真空光速的严格限制。

\subsection{3D集成}
\label{sec:cpu:3D Integration}

三维集成（3DI）是将非常薄的硅芯片垂直堆叠粘合在一起的工艺。
这项工艺带来了一些潜在的好处，但也带来了
重大的制造挑战~\cite{JohnKnickerbocker2008:3DI}。

\begin{figure}
\centering
\resizebox{3in}{!}{\includegraphics{cpu/3DI}}
\caption{3D集成的延迟优势}
\label{fig:cpu:Latency Benefit of 3D Integration}
\end{figure}

也许3DI所带来的最大好处是降低了系统的路径长度，
如\cref{fig:cpu:Latency Benefit of 3D Integration}所示。
一个3厘米的硅芯片可以被四块重叠的1.5厘米芯片替代，
要知道每一层都非常薄。
不仅如此，如果能小心地设计和加工，
长的水平方向电连接（既慢又耗电）可以被短的垂直方向电连接所取代，
后者既更快又更节能。

如果你不能让光跑得更快，就把你的器件做得更小！

不过，时钟逻辑结构级别造成的延迟是无法通过3D集成来降低的，
并且诸如生产、测试、供电和散热等3D集成中的重大问题
必须得到解决才有望产品化。
散热问题可能需要用基于金刚石的半导体来解决，
金刚石是热的良好导体，但却是电的绝缘体。
据说生长大尺寸的单晶金刚石仍然很困难，
更别说将金刚石切割成晶圆了。
不仅如此，上述这些技术看起来不大可能像某些人习惯的那样让性能出现指数级增长。
不过，它们可能是通往已故Jim Gray所说的
``冒烟的毛茸茸的高尔夫球''~\cite{JimGray2002SmokingHairyGolfBalls}道路上的必要步骤。

\begin{figure}
\centering
\resizebox{3in}{!}{\includegraphics{cpu/High_Bandwidth_Memory_schematic}}
\caption{3D集成示例}
\ContributedBy{fig:cpu:Example 3D Integration}{Shmuel Csaba Otto Traian, CCSA4.0}
\end{figure}

与此同时，大约2020年，chiplet\footnote{
	\url{https://en.wikipedia.org/wiki/Chiplet}}
技术的进步在这个方向上取得了重大进展，
如\cref{fig:cpu:Example 3D Integration}所示。
一些研究人员正在更进一步，在单个芯片内堆叠
晶体管~\cite{SamuelKMoore2020StackedCMOS,MarkoRadosavljevic2022StackedCMOS}。

\subsection{专用加速器}
\label{sec:cpu:Special-Purpose Accelerators}

用通用CPU来处理某个专门问题，通常会将大量的时间和能源消耗
在和问题无关的地方。
例如，当对一对向量进行点积操作时，通用CPU一般会使用一个循环
（通常已展开）和一个循环计数器。
对指令解码、递增循环计数器、测试计数器以及跳转回循环顶部，
从某种意义上来说都是无用功：
真正的目标是计算两个向量对应元素的乘积。
因此，在硬件上设计一个专用于向量乘法的部件
可以让这项工作完成得既快速又节能。

这就是在许多商用微处理器上出现向量指令的原因。
这些指令可以同时操作多个数据，
能让点积计算减少指令解码和循环开销的消耗。

类似地，专用硬件可以更有效地进行加密/解密、
压缩/解压缩、编码/解码等许多任务。
不过不幸的是，这种效率提升并非免费的。
包含专用硬件的计算机系统将需要更多的晶体管，
即使在不用时也会带来能源消耗。
软件也必须进行修改以利用专用硬件的长处，
同时这种专用硬件又必须有足够广泛的通用性，
这样高昂的前期设计费用才能摊到足够多的用户身上，
让专用硬件的价格变得可以承受。
部分由于以上经济考虑因素，专用硬件迄今只出现在少数
应用领域，包括图形处理（GPU）、向量处理器（MMX、SSE
和VMX指令），以及在较小程度上的加密和压缩。
即使在这些领域，实现预期的性能提升也并非总是容易的，
例如，由于热节流~\cite{VladKrasnov2017SIMDfreqscale,DanielLemire2018SIMDfreqscale,TravisDowns2020SIMDfreqscale}。

与服务器和PC领域不同，智能手机长期以来使用了
各种各样的硬件加速器。
这些硬件加速器通常用于媒体解码，
以至于一个高端MP3播放器可能能够播放数分钟的音频——
其CPU在整个过程中完全断电。
这些加速器的目的是提高能源效率从而延长电池寿命：
专用硬件通常可以比通用CPU更高效地进行计算。
这是\cref{sec:intro:Generality}中指出的原则的又一个例子：
\IX{Generality}几乎从来都不是免费的。

尽管如此，鉴于\IXaltr{Moore's-Law}{Moore's Law}驱动的
单线程性能提升的终结，可以安全地假设将会出现越来越多种类的专用硬件。
例如，在2020年代中期，许多人押注于专用于人工智能和
机器学习工作负载的加速器。

\subsection{现有的并行软件}
\label{sec:cpu:Existing Parallel Software}

虽然多核CPU曾经让计算行业惊讶了一把，
但事实上基于共享内存的并行计算机至少从1980年代起
就已经随时可用（尽管价格不菲）。
这段时间足以让一些重要的并行软件登上舞台，
事实也确实如此。
并行操作系统已经非常常见，并行线程库、
并行关系数据库管理系统和并行数值软件也是如此。
这些现有的并行软件在解决我们可能遇到的并行难题上已经走了很长一段路。

也许最常见的例子是并行关系数据库管理系统。
与单线程程序不同，并行关系数据库管理系统一般用高级脚本语言书写，
并发地访问位于中心的关系数据库。
在这种高度并行化的系统中，只有数据库本身需要直接处理并行性。
当这种方式奏效时，真是一个非常好的技巧！
